\section{Experiments}
\label{sec:experiments}

We evaluate \ours{} along four lines of evidence. First, we show a
headline SOTA on FACED EEG emotion classification
(\cref{sec:faced}). Second, we demonstrate cross-modality generality on
six image and text benchmarks (\cref{sec:image}), from which emerges a
predictive scaling law
(\cref{sec:scaling}). Third, we compare head-to-head against SWA,
snapshot, and 10-seed deep ensembles on accuracy \emph{and} calibration
(\cref{sec:baselines,sec:calib}). Fourth, we open the mechanism: a
theoretical decomposition on FACED, weight-space PCA and
mode-connectivity probes, and an endpoint-ratio sweep that answers
``why $T_B=150$?'' empirically
(\cref{sec:decomp,sec:curves,sec:endpoint_sweep,sec:basin}).

\subsection{Setup}

We test \ours{} on 9 benchmarks spanning three modalities (EEG, images,
text) and multiple architectures:

\begin{itemize}[leftmargin=*]
\item \textbf{FACED} \cite{liu2024faced}: 9-class EEG emotion classification,
10{,}332 samples (6{,}720 train / 1{,}680 val / 1{,}932 test).
Architecture: EMOD axial-transformer with depth=6 (3 pretrained +
3 random init), $n_\mathrm{filters}=128$.
\item \textbf{SEED-V} \cite{zheng2019investigating}: 5-class EEG emotion
classification. Same EMOD d6 architecture.
\item \textbf{CIFAR-10}, \textbf{CIFAR-100} \cite{krizhevsky2009learning}:
image classification with a small CNN (64$\to$128$\to$128$\to$256 channels,
$\sim$500K params).
\item \textbf{MNIST}, \textbf{Fashion-MNIST}: with small MLP (3 layers,
FC-FC-FC) and small CNN respectively.
\item \textbf{SVHN}: house number recognition with small CNN.
\end{itemize}

For each dataset, we train 5 seeds with $T_A$ epochs and 5 seeds with
$T_B$ epochs (same augmentation, same optimizer, same LR schedule except
for $T_\mathrm{max}$), then compare:
(a)~individual-seed accuracy,
(b)~5-seed same-length deep ensemble accuracy (gold-standard baseline),
(c)~10-model mixed-length \ours{} ensemble accuracy.

\subsection{EEG Emotion Classification: A New FACED SOTA}
\label{sec:faced}

\begin{table}[h]
\centering
\caption{FACED 9-class EEG emotion classification. EMOD d6 axial transformer,
aug+LS+rand9 KD recipe. Our 10-model mixed-length ensemble is a new SOTA,
$+5.3$~pp over the EMOD AAAI 2026 baseline.}
\label{tab:faced}
\begin{tabular}{lcc}
\toprule
Method & Balanced Acc & $\Delta$ over baseline \\
\midrule
REVE NeurIPS 2025 \cite{kostas2025reve} & 0.500 & $-$0.142 \\
EMOD AAAI 2026 \cite{chen2026emod} & 0.642 & --- \\
EMOD d3 + aug + LS=0.1 + rand9 KD (ours, cycle 72) & 0.6467 $\pm$ 0.007 & $+$0.005 \\
EMOD d6 (3 pretrained + 3 random) single model & 0.6581 $\pm$ 0.007 & $+$0.016 \\
EMOD d6 5-seed ensemble & 0.6798 & $+$0.038 \\
EMOD d6 10-seed same-length ensemble & 0.6838 & $+$0.042 \\
\rowcolor{blond}
\textbf{EMOD d6 + d6\_e150 10-model \ours{}} & \textbf{\facedacc{}} & \textbf{$+$0.053} \\
\bottomrule
\end{tabular}
\end{table}

\begin{table}[h]
\centering
\caption{SEED-V 5-class EEG emotion classification. EMOD d6 cross-dataset
transfer. No 150ep checkpoints available, so we only report the standard
5-seed ensemble gain.}
\label{tab:seedv}
\begin{tabular}{lc}
\toprule
Method & Balanced Acc \\
\midrule
EMOD d3 + aug + KD (ours, cycle 244) & 0.3744 \\
EMOD d6 + aug + rand9 KD, individual mean & 0.3743 $\pm$ 0.006 \\
\rowcolor{blond}
\textbf{EMOD d6 5-seed ensemble} & \textbf{\seedvacc{}} (\seedvgain{}) \\
\bottomrule
\end{tabular}
\end{table}

Crucially, \ours{} beats the 10-seed same-length deep ensemble
($0.6838$) by $+1.1$~pp---a margin larger than the gap between the
5-seed and 10-seed same-length ensembles---so the gain is not a matter
of buying more seeds. Two trajectory-length groups at
$N{=}5$ each outperform $N{=}10$ seed diversity. SEED-V confirms the
picture at $+3.4\%$ cross-dataset.

\subsection{Cross-Modality Validation: Images and Text}
\label{sec:image}

Having established the EEG headline, we next ask whether the same
single-integer change transports to other modalities and architectures.
We use standard small-CNN, MLP, ResNet-18, and DistilBERT backbones on
seven additional benchmarks whose individual-model accuracies span the
range from $\sim\!0.69$ (CIFAR-100) to $\sim\!0.98$ (MNIST).

\begin{table}[h]
\centering
\caption{\ours{} on four image datasets with small CNN/MLP architectures.
Each row: 5 seeds × 50ep (group A) + 5 seeds × 75ep (group B), mixed 10-model
ensemble. Gain scales inversely with individual accuracy.}
\label{tab:image}
\begin{tabular}{lcccc}
\toprule
Dataset & Arch & Ind.\ mean & 5-seed ensemble & \textbf{10-model mixed} \\
\midrule
CIFAR-10 & Small CNN & 0.9015 & 0.9209 (+0.0194) & \textbf{\ciftenacc{} (+0.0238)} \\
CIFAR-10 & ResNet-18 & 0.9452 & 0.9563 (+0.0111) & \textbf{0.9599 (+0.0147)} \\
CIFAR-100 & Small CNN & 0.6886 & 0.7225 (+0.0339) & \textbf{0.7353 (+0.0467)} \\
CIFAR-100 & ResNet-18 & 0.7707 & 0.8020 (+0.0313) & \textbf{0.8104 (+0.0397)} \\
MNIST & MLP & 0.9849 & 0.9859 (+0.0010) & \textbf{0.9863 (+0.0014)} \\
Fashion-MNIST & Small CNN & 0.9313 & 0.9418 (+0.0105) & \textbf{0.9434 (+0.0121)} \\
SST-2 & DistilBERT & 0.9062 & 0.9094 (+0.0032) & 0.9094 (+0.0032) \\
SVHN$^{\dagger}$ & Small CNN & -- & -- & -- \\
\bottomrule
\end{tabular}
\\[0.4em]
{\footnotesize $\dagger$: SVHN run excluded from main table. The short-CNN
trainer used for this sweep had a training instability that caused 7 of
10 seeds to fail to escape the uniform-prediction basin. A clean SVHN
rerun is scheduled (see supplementary).}
\end{table}

\subsection{A Predictive Scaling Law for the Gain}
\label{sec:scaling}

A striking pattern emerges from the nine benchmarks together: the
mixed-length ensemble gain is a \emph{predictable function} of how far
each individual model is from ceiling. Every added modality confirms
the same fit.
A least-squares linear fit of gain on headroom gives slope $0.057$,
intercept $0.011$, Pearson $r=0.753$, $p=0.019$, $R^2=0.57$; a power-law
fit $\text{gain} = c \cdot \text{headroom}^{\alpha}$ gives exponent
$\alpha = 0.85$ (log-space $R^2 = 0.72$). Both fits support the intuition
that the gain is governed by the fraction of samples still correctable
by ensembling.
On CIFAR-100, we note that the mixed ensemble at $0.7353$ is slightly
below the best single-length 5-seed ensemble at $B=75$~epochs
($0.7367$); the mixed gain is positive when compared against the average
of the two same-length ensembles (\emph{mean}($A$-5-seed, $B$-5-seed)
$=0.7296$, mixed $-$ mean $=+0.0057$). CIFAR-100 is the one dataset
where the longer schedule $B$ alone beats the mixed ensemble by a narrow
margin; the mixed ensemble still wins over the matched same-length
baseline.

\begin{figure}[tb]
\centering
\includegraphics[width=0.75\textwidth]{fig_scaling_law}
\caption{Ensemble gain vs.\ headroom $(1 - \text{individual accuracy})$
across nine benchmarks. Dashed red: through-the-origin fit
$\text{gain} = 0.086 \cdot \text{headroom}$. Dotted green: ordinary
least-squares fit with intercept ($r=0.753$, $p=0.019$, $R^2=0.57$).
A power-law fit gives exponent $\alpha=0.85$.}
\label{fig:scaling}
\end{figure}

\begin{table}[h]
\centering
\caption{Ensemble gain scales inversely with individual accuracy.}
\label{tab:scaling}
\begin{tabular}{lccc}
\toprule
Dataset & Ind.\ accuracy & Mixed gain & Ratio: gain/(1-ind.) \\
\midrule
FACED 9c EEG           & 0.6581 & +0.0367 & 0.107 \\
SEED-V 5c EEG          & 0.3743 & +0.0340 & 0.054 \\
CIFAR-100 CNN          & 0.6886 & +0.0467 & 0.150 \\
CIFAR-100 ResNet-18    & 0.7707 & +0.0397 & 0.173 \\
CIFAR-10 CNN           & 0.9015 & +0.0238 & 0.242 \\
CIFAR-10 ResNet-18     & 0.9452 & +0.0147 & 0.268 \\
Fashion-MNIST CNN      & 0.9313 & +0.0121 & 0.176 \\
SST-2 DistilBERT       & 0.9062 & +0.0032 & 0.034 \\
MNIST MLP              & 0.9849 & +0.0014 & 0.093 \\
\bottomrule
\end{tabular}
\end{table}

\subsection{Head-to-Head Against SWA, Snapshot, and Deep Ensembles}
\label{sec:baselines}

The scaling law predicts \emph{when} \ours{} helps; we next ask
\emph{how much} it helps against standard alternatives at matched
compute. On CIFAR-10 with the same small CNN we compare
(i)~a 10-seed same-length deep ensemble, (ii)~Stochastic Weight
Averaging (SWA start at epoch 30 of a 60-epoch run)~\cite{izmailov2018averaging},
and (iii)~a snapshot ensemble with 5 cyclic-LR snapshots from a single
run~\cite{huang2017snapshot}. All methods share hyperparameters and
training code; \ours{} uses roughly the same compute as the 10-seed
deep ensemble ($5\times50 + 5\times75 = 625$ vs.\ $10\times50 = 500$
epochs, a $1.25\times$ increase).

\begin{table}[h]
\centering
\caption{Head-to-head comparison with SWA, snapshot ensembles, and standard
10-seed deep ensembling on CIFAR-10 with the same small CNN.
Same total compute budget.}
\label{tab:cifar10_baselines}
\begin{tabular}{lcc}
\toprule
Method & Test Accuracy \\
\midrule
Single model, 50 epochs, mean of 10 seeds & 0.8992 \\
Standard 5-seed deep ensemble (50ep) & 0.9209 (+0.0217) \\
Standard 10-seed deep ensemble (50ep) & 0.9223 (+0.0231) \\
SWA (60ep with SWA starting at 30) & 0.8918 ($-0.0074$) \\
Snapshot ensemble (1 run, 5 cyclic-LR snapshots) & 0.8946 ($-0.0046$) \\
\rowcolor{blond}
\textbf{\ours{} (5$\times$50ep + 5$\times$75ep, 10 models mixed)} & \textbf{\ciftenacc{} (+0.0261)} \\
\bottomrule
\end{tabular}
\end{table}

\begin{figure}[tb]
\centering
\includegraphics[width=0.85\textwidth]{cifar10_baselines_bar}
\caption{CIFAR-10 head-to-head ensemble comparison at matched compute.
\ours{} beats every baseline including the 10-seed standard deep
ensemble. SWA and snapshot ensembling both \emph{hurt} relative to the
single-model mean on this small-CNN trainer, likely because they
average over a narrow region of weight space.}
\label{fig:cifar10_bar}
\end{figure}

\subsection{Calibration: NLL, ECE, and Brier at Matched Compute}
\label{sec:calib}

Ensembles are evaluated on calibration as much as
accuracy~\cite{ashukha2020pitfalls}. We rerun the CIFAR-10 head-to-head
on fresh seeds and save full softmax logits, then report accuracy, NLL
(negative log-likelihood, nats), ECE (15-bin expected calibration
error), and Brier score. (Accuracies differ slightly from
\cref{tab:cifar10_baselines} because of the fresh seeds.)

\begin{table}[h]
\centering
\caption{CIFAR-10 full-calibration comparison (same small CNN, fresh
seeds; accuracy numbers differ slightly from Table~\ref{tab:cifar10_baselines}
because of the fresh seeds). \ours{} sweeps all four metrics over SWA,
snapshot, and a 10-seed same-length deep ensemble. Lower is better for
NLL, ECE, Brier.}
\label{tab:cifar10_calib}
\begin{tabular}{lcccc}
\toprule
Method & Acc $\uparrow$ & NLL $\downarrow$ & ECE $\downarrow$ & Brier $\downarrow$ \\
\midrule
Single model (50ep) mean of 5 seeds & 0.8904 & 0.3469 & 0.0328 & 0.1637 \\
SWA (60ep)                          & 0.8865 & 0.3537 & 0.0275 & 0.1688 \\
Snapshot (5 cyclic snaps)           & 0.8882 & 0.3372 & 0.0180 & 0.1638 \\
5-seed ens 50ep                     & 0.9058 & 0.2813 & 0.0126 & 0.1389 \\
5-seed ens 75ep                     & 0.9146 & 0.2557 & 0.0104 & 0.1253 \\
5-seed ens 100ep                    & 0.9190 & 0.2479 & 0.0125 & 0.1196 \\
10-seed same-length 50ep            & 0.9078 & 0.2722 & 0.0101 & 0.1355 \\
\rowcolor{blond}
\textbf{\ours{} (5$\times$50ep + 5$\times$65ep)}   & 0.9126 & 0.2607 & 0.0098 & 0.1290 \\
\rowcolor{blond}
\textbf{\ours{} (5$\times$50ep + 5$\times$75ep)}   & 0.9149 & 0.2566 & 0.0107 & 0.1277 \\
\rowcolor{blond}
\textbf{\ours{} (5$\times$50ep + 5$\times$100ep)}  & \textbf{0.9168} & \textbf{0.2489} & \textbf{0.0096} & \textbf{0.1236} \\
\bottomrule
\end{tabular}
\end{table}

\paragraph{Findings.}
(i)~\ours{} with $T_B=100$ attains the best NLL, the best ECE, and the
best Brier across the whole table. (ii)~SWA \emph{hurts} NLL relative
to a single model ($0.3537$ vs.\ $0.3469$): weight averaging collapses
the predictive distribution into a point estimate, while \ours{}
preserves it by averaging outputs. (iii)~Across the endpoint-ratio
sweep ($T_B \in \{65, 75, 100\}$) all four calibration metrics improve
monotonically with $T_B$, consistent with a longer tail visiting more
of the basin.

\subsection{Mechanism (I): Theoretical Decomposition of the Gain}
\label{sec:decomp}

Having shown the gain across tasks and against baselines, we now open
the mechanism box. On FACED d6 we save all 10 checkpoints and decompose
the gain into within-group and cross-group contributions.

\begin{table}[h]
\centering
\caption{Ensemble gain decomposition on FACED d6. 100ep and 150ep
individual-seed means are identical within 0.0001; the 150ep ensemble is
stronger because 150ep seeds are more diverse.}
\label{tab:decomposition}
\begin{tabular}{lc}
\toprule
Ensemble & Balanced Acc \\
\midrule
Individual 100ep mean (5 seeds) & 0.6581 \\
Individual 150ep mean (5 seeds) & 0.6581 \\
\midrule
5-seed 100ep ensemble & 0.6798 ($+0.0218$) \\
\textbf{5-seed 150ep ensemble} & \textbf{0.6919 ($+0.0338$)} \\
\textbf{10-model mixed} & \textbf{\facedacc{} ($+0.0367$)} \\
\midrule
Mixed gain over 100ep alone & $+0.0150$ \\
\bottomrule
\end{tabular}
\end{table}

\paragraph{Decision-boundary analysis.}
Of the 1{,}932 FACED test samples, 50.9\% have unanimous prediction from
all 10 models; 13.1\% (253 samples) have a \emph{different} group
consensus between 100ep and 150ep. On those 253 samples, the 100ep
consensus accuracy is 28.85\%, the 150ep consensus accuracy is 29.25\%,
and the mixed ensemble accuracy is 39.53\%---a gain of $+10.3$ percentage
points on a subset of 13.1\% of the test set. The predicted overall gain
$10.3 \times 0.131 = 1.35$~pp closely matches the observed $+1.50$~pp.

\subsection{Mechanism (II): Same Seed, Different $T_\mathrm{max}$, Different Basins}
\label{sec:curves}

The decomposition predicts the magnitude of the gain; the next question
is whether the two groups really occupy \emph{different local minima} or
just two points in one wider basin. Following
Fort~\emph{et al.}~\cite{fort2019deepensembles}, we train two CIFAR-10
small CNNs with the \emph{same} random seed (2025) at $T_A=50$ and
$T_B=75$ epochs and linearly interpolate between their weights
$w(\alpha)=(1-\alpha) w_A + \alpha w_B$, $\alpha \in [0,1]$, measuring
the train and test loss at each step.

\begin{figure}[tb]
\centering
\includegraphics[width=0.85\textwidth]{loss_landscape_interp}
\caption{Train/test loss along the linear interpolation between a
$T_A=50$ and a $T_B=75$ CIFAR-10 checkpoint (same seed). Both endpoints
are low-loss minima (train $\approx 0.05$--$0.12$), but the midway
linear path has a \textbf{+0.25 nats train-loss barrier}
and accuracy drops to $0.843$ ($-6$~pp). The two endpoints are not
in the same basin --- they are separated by a loss ridge, directly
confirming that trajectory-length variation drifts weights into a
neighboring local minimum.}
\label{fig:landscape}
\end{figure}

\paragraph{Result.} The train-loss barrier is $+0.2497$ and the
test-loss barrier is $+0.1266$. The test accuracy at $\alpha=0.5$ is
$0.8434$, compared to $0.8905$ at $\alpha=0$ and $0.9018$ at
$\alpha=1$ --- a $-6$~pp drop. This is a textbook ``different basins''
signature. Same seed, different $T_\mathrm{max}$, different minima.

\subsection{Mechanism (III): Same Final Accuracy, Different Trajectories}

If the two groups reach different minima, the per-epoch validation
curves should reveal the divergence. We plot the full training curves
for all 10 FACED d6 seeds. The two trajectories converge to the
\emph{same} final accuracy but follow different routes to get
there---the 150ep group is still at $\approx$0.25 of initial LR at
epoch 100 and spends 50 more epochs drifting before stopping.

\begin{figure}[tb]
\centering
\includegraphics[width=0.75\textwidth]{training_curves_100_vs_150}
\caption{Validation BAcc curves for 100ep and 150ep d6 seeds on FACED. The
100ep trajectory (solid) converges around epoch 100 at BAcc $\approx$0.677.
The 150ep trajectory (dashed) is slightly behind at epoch 100 (still at
BAcc $\approx$0.672, with LR $=0.25\times$ initial), then slowly ascends
to the same final BAcc by epoch 149. The extra 50 epochs add zero mean
accuracy but substantial parameter drift.}
\label{fig:curves}
\end{figure}

\begin{table}[h]
\centering
\caption{Per-epoch val BACC for 100ep and 150ep d6 FACED seeds (mean over
5 seeds). Identical final accuracy, different trajectories.}
\label{tab:trajectories}
\begin{tabular}{lcc}
\toprule
Epoch & 100ep group BAcc & 150ep group BAcc \\
\midrule
25 & 0.6267 $\pm$ 0.0118 & 0.6186 $\pm$ 0.0057 \\
50 & 0.6513 $\pm$ 0.0090 & 0.6516 $\pm$ 0.0046 \\
75 & 0.6718 $\pm$ 0.0046 & 0.6577 $\pm$ 0.0079 \\
99 & 0.6771 $\pm$ 0.0059 & 0.6722 $\pm$ 0.0058 \\
149 & --- & 0.6776 $\pm$ 0.0040 \\
\bottomrule
\end{tabular}
\end{table}

\subsection{Ablation: The Endpoint Ratio is a Robust Operating Regime, Not a Tuned Constant}
\label{sec:endpoint_sweep}

A natural concern is that $T_B=150$ might be a cherry-picked sweet
spot. To rule this out, we run the full sweep
$T_B \in \{110, 120, 130, 140, 150, 200\}$ on FACED EMOD d6 with $T_A=100$
held fixed. \cref{tab:endpoint_sweep} and \cref{fig:endpoint_sweep}
report, for each $T_B$, the mean individual BAcc of the $T_B$ group,
the $T_B$-only 5-seed ensemble, and the mixed 10-model ensemble formed
by $T_A=100 \cup T_B$.

\begin{table}[h]
\centering
\caption{FACED EMOD d6 endpoint-ratio ablation. All rows use the
same $T_A=100$ group (mean ind.\ 0.6581, 5-seed ens.\ 0.6798). Mixed
10-model BAcc peaks at $T_B=150$ (our default) and degrades both for
smaller drifts ($T_B \le 140$) and for over-long drifts
($T_B=200$). Source: \texttt{llm\_steering/diag\_p2\_d4\_endpoint\_T*.json}.}
\label{tab:endpoint_sweep}
\begin{tabular}{cccc}
\toprule
$T_B$ & Mean ind.\ $T_B$ & 5-seed ens.\ $T_B$ & \textbf{Mixed 10-model} \\
\midrule
100 (=$T_A$) & 0.6581 & 0.6798 & 0.6798 \\
110 & 0.6612 & 0.6809 & 0.6842 \\
120 & 0.6418 & 0.6740 & 0.6805 \\
130 & 0.6308 & 0.6758 & 0.6778 \\
140 & 0.6292 & 0.6730 & 0.6750 \\
\rowcolor{blond}
\textbf{150} & \textbf{0.6581} & \textbf{0.6919} & \textbf{0.6948} \\
200 & 0.5960 & 0.6425 & 0.6633 \\
\bottomrule
\end{tabular}
\end{table}

\begin{figure}[tb]
\centering
\includegraphics[width=0.72\textwidth]{fig_endpoint_sweep}
\caption{Mixed 10-model BAcc on FACED as a function of $T_B$ with
$T_A=100$ fixed. Dashed green: 100ep-only 5-seed ensemble baseline.
The endpoint ratio $\rho = T_B/T_A = 1.5$ is a genuine operating-regime
sweet spot, consistent with the PCA-basin and mode-connectivity
analyses (\cref{sec:basin}).}
\label{fig:endpoint_sweep}
\end{figure}

\paragraph{Interpretation.}
The sweep traces an interpretable curve that matches the basin
geometry. For intermediate $\rho \in (1.1, 1.4]$ the $T_B$ trajectory
has not yet spent enough low-LR tail time to reach a distinct basin,
and its individual accuracy is also transiently depressed because the
cosine tail is stretched; the mixed ensemble sits just above the
$T_A$-only baseline. At $\rho = 1.5$ the mean individual BAcc recovers
to $0.6581$ (within $0.0001$ of $T_A$) while the run has spent 50
additional epochs in the productive tail and settles in a
geometrically distinct basin (\cref{sec:basin}); this is exactly where
the mixed gain is maximal ($+0.015$ over $T_A$-only). At $\rho = 2.0$
the long tail leaves the productive regime, individual accuracy drops
to $0.596$, and the mixed ensemble is dragged below the baseline. The
takeaway: \ours{} operates in a broad regime $\rho \in [1.4, 1.6]$
with $\rho = 1.5$ as a task-robust default and $\rho \ge 2.0$ to be
avoided.

\subsection{Mechanism (IV): Different Basins, and Why Weight-Space Averaging Fails}
\label{sec:basin}

The endpoint sweep's sweet spot at $\rho=1.5$ coincides with the regime
where the two groups reach distinct basins. We verify basin separation
directly in two complementary ways on FACED: a weight-space PCA of the
10 checkpoints, and a linear mode-connectivity probe between cross-group
pairs.

\paragraph{Weight-space PCA.}
Flattening the 10 FACED d6 checkpoints (5 seeds $\times$ 2 endpoints) into
a matrix and taking its first three principal components shows that PC1
captures $>\!99\%$ of the between-checkpoint variance and perfectly
separates the two groups (\cref{fig:pca_basin}). Concretely, the mean
pairwise distance between the two group centroids is $5166$ and the mean
pairwise distance within each group is $1587$, a between/within ratio of
$\mathbf{3.25}$. The 100ep and 150ep runs are thus geometrically
distinct in weight space: they inhabit two separate
neighborhoods, not one broad basin visited at two points along a cosine
schedule.

\begin{figure}[tb]
\centering
\includegraphics[width=0.72\textwidth]{fig_pca_basin}
\caption{Weight-space PCA of all 10 FACED d6 checkpoints. PC1 accounts
for $>\!99.9\%$ of variance and fully separates the 100ep and 150ep
groups. Between/within ratio $= 3.25$. Source:
\texttt{reports/p2\_pca\_basin.json}.}
\label{fig:pca_basin}
\end{figure}

\paragraph{Parameter-space averaging fails; prediction-space averaging works.}
To check whether the two groups' minima are connected by a low-loss
linear path (same basin) or separated by a ridge (different basins), we
linearly interpolate between cross-group pairs
$\theta(\lambda) = (1-\lambda)\theta_{100} + \lambda\theta_{150}$ for
$\lambda \in \{0, 0.1, \ldots, 1\}$ and evaluate val and test loss at
each $\lambda$ (\cref{fig:mode_conn}, red curve). The interpolated
weights produce \emph{near-chance} performance for every
$\lambda$ (test loss $\approx 2.18$ nats, test BAcc $0.24$--$0.26$ vs.\
$0.66$ at the endpoints): the linear path between a 100ep and a 150ep
minimum is flat at chance, indicating that the two endpoints are not
even connected by a low-loss straight line in weight space. For reference,
within-group interpolation ($\lambda$-path between two 100ep seeds or
two 150ep seeds, blue/orange dashed) traces the usual mode-connectivity
barrier: test loss peaks near $\lambda=0.5$ at $\approx 2.0$ and
recovers to the endpoint value.

This is a crisp design signal: \textbf{prediction-space averaging
succeeds exactly where parameter-space averaging fails}. Output
averaging is agnostic to the geometry of the weight space and only
requires that the two predictive distributions disagree on the right
samples, which is what our decomposition already confirmed. It also
explains why SWA (which averages weights) \emph{hurts} NLL on this
trainer (\cref{sec:calib}): weight averaging across trajectories that
live in different basins collapses the ensemble's uncertainty rather
than combining it.

\begin{figure}[tb]
\centering
\includegraphics[width=0.95\textwidth]{fig_mode_connectivity}
\caption{Linear interpolation loss on FACED. Solid red:
$\theta_{100}\!\rightarrow\!\theta_{150}$ (cross-group, different seeds);
loss stays at chance level ($\approx 2.18$ nats) along the entire path,
confirming the two groups live in different basins. Dashed blue/orange:
within-group 100ep-to-100ep and 150ep-to-150ep interpolations; these
show the expected $\sim\!1.0$-nat mode-connectivity barrier. Takeaway:
weight averaging between groups is not a meaningful ensemble
substitute; prediction-space averaging preserves the gain.
Source: \texttt{reports/p2\_mode\_connectivity.json}.}
\label{fig:mode_conn}
\end{figure}

\paragraph{Disagreement decomposition (visual).}
\cref{fig:disagree} shows the accuracy of each estimator on the 253
disagreement samples (13.1\% of FACED test). The mixed 10-model
ensemble (rightmost, yellow) achieves $39.53\%$ accuracy on this
subset, $+10.68$~pp above the 100ep consensus and $+10.28$~pp above the
150ep consensus; the gain is cleanly localized to the disagreement set.

\begin{figure}[tb]
\centering
\includegraphics[width=0.78\textwidth]{fig_disagreement_decomposition}
\caption{Per-method accuracy on the 253 disagreement samples
(13.1\% of FACED test). The two group consensuses are nearly tied
($28.85\%$ vs.\ $29.25\%$); each group's 5-seed ensemble recovers a few
percentage points; the 10-model mixed \ours{} ensemble recovers the
most. Source: \texttt{reports/p2\_disagreement\_deep\_dive.json}.}
\label{fig:disagree}
\end{figure}
